
\documentclass[11pt,letterpaper]{article}

\usepackage[margin=1in]{geometry}
\usepackage[T1]{fontenc}
\usepackage[utf8]{inputenc}
\usepackage{lmodern}
\usepackage{microtype}
\usepackage{amsmath,amssymb,amsthm,mathtools}
\usepackage{booktabs}
\usepackage{tabularx}
\usepackage{array}
\usepackage{enumitem}
\usepackage{xcolor}
\usepackage{listings}
\usepackage{longtable}
\usepackage{graphicx}
\usepackage{hyperref}
\usepackage[nameinlink,noabbrev]{cleveref}

\definecolor{codegray}{gray}{0.95}
\definecolor{linkblue}{RGB}{25,75,125}
\hypersetup{
  colorlinks=true,
  linkcolor=linkblue,
  citecolor=linkblue,
  urlcolor=linkblue,
  pdftitle={Link Chat: A Reproducible Protocol Core for Endpoint-Local State Evolution},
  pdfauthor={Link Chat Research Project}
}

\lstset{
  basicstyle=\ttfamily\small,
  backgroundcolor=\color{codegray},
  frame=single,
  rulecolor=\color{black!20},
  columns=fullflexible,
  keepspaces=true,
  breaklines=true,
  showstringspaces=false,
  captionpos=b
}

\newcolumntype{Y}{>{\raggedright\arraybackslash}X}
\newcolumntype{P}[1]{>{\raggedright\arraybackslash}p{#1}}

\newtheorem{definition}{Definition}[section]
\newtheorem{theorem}[definition]{Theorem}
\newtheorem{lemma}[definition]{Lemma}
\newtheorem{proposition}[definition]{Proposition}
\newtheorem{corollary}[definition]{Corollary}
\newtheorem{remark}[definition]{Remark}

\newcommand{\Adv}{\mathsf{Adv}}
\newcommand{\negl}{\mathsf{negl}}
\newcommand{\PPT}{\mathsf{PPT}}
\newcommand{\View}{\mathsf{View}}
\newcommand{\State}{\mathsf{State}}
\newcommand{\Package}{\mathsf{Package}}
\newcommand{\Turn}{\mathsf{Turn}}
\newcommand{\Enc}{\mathsf{Enc}}
\newcommand{\Dec}{\mathsf{Dec}}
\newcommand{\Seal}{\mathsf{Seal}}
\newcommand{\Open}{\mathsf{Open}}
\newcommand{\Hash}{\mathsf{Hash}}
\newcommand{\Sign}{\mathsf{Sign}}
\newcommand{\Verify}{\mathsf{Verify}}
\newcommand{\KDF}{\mathsf{KDF}}
\newcommand{\sample}{\xleftarrow{\$}}
\newcommand{\code}[1]{\texttt{#1}}

\title{Link Chat: A Reproducible Protocol Core for\
Endpoint-Local State Evolution}
\author{Link Chat Research Project}
\date{September 6, 2026}

\begin{document}

\maketitle

\begin{abstract}
This paper presents Link Chat, a small two-party protocol core based on the Link Chat v1.0
research specification \cite{linkchat-spec}, designed to study how
standard cryptographic mechanisms interact with endpoint-local state evolution. The
construction is intentionally not introduced as a new cipher or as a complete messaging
service. Its central design combines independent receiver secret state at each endpoint,
strictly alternating application turns, authenticated public receiver packages, and an
atomic package rotation performed only after an authenticated application input has been
fully validated. Transport, mailbox, persistence, and application components are separated
from the authority to advance the application state.

The paper gives a complete protocol description, a canonical byte-level encoding, a fixed
4096-byte envelope format, typed primitive interfaces, endpoint algorithms, failure rules,
storage commit semantics, and reproducible test procedures. It also formalizes a complete
adaptive adversary view containing protocol outputs and metadata such as timing, logical
length, envelope length, acknowledgements, errors, and state-update counts. We state the
structural properties checked by the associated Lean development and formulate a
conditional computational theorem for a supplied concrete probability model, PPT
adversary witness, wrapper-game certificates, and hybrid reduction bound. The theorem is
deliberately conditional: the Lean development does not claim to derive the computational
security of X25519, ML-KEM-768, HKDF, ChaCha20-Poly1305, Ed25519, or an operating-system
CSPRNG from their implementations.
\end{abstract}

\noindent\textbf{Keywords:} protocol state machines, receiver-state rotation, non-interference,
post-compromise reasoning, game-based security, Lean, reproducible wire formats.

\tableofcontents

\section{Introduction}
\label{sec:introduction}

Modern end-to-end messaging systems combine cryptographic primitives with stateful
protocol logic. A protocol can use well-studied primitives while still leaving important
questions unanswered about how malformed, replayed, delayed, or adversarially selected
inputs affect honest local state. The protocol specification isolates
this question.

The core design question is the following:

\begin{quote}
Can a peer-controlled application input influence authenticated message processing without
acquiring authority to choose the honest endpoint's next private receiver state?
\end{quote}

The answer pursued by Link Chat is architectural rather than algorithmic. Each endpoint
owns only its local receiver secret state. The peer receives a public projection of that
state, consisting of public keys, a key identifier, a mailbox token, a generation value,
expiration information, a chain link, and an identity-bound authentication value. The
application protocol advances in a strict alternating order. A valid receive operation
does not mutate the current state in place: it constructs a complete candidate state,
persists it through an explicit prepare/commit boundary, and exposes the application
plaintext only after successful commit.

The result is a small protocol core with four intended uses:
\begin{enumerate}[leftmargin=*]
  \item a precise object and state-transition specification for independent
        implementations;
  \item a test target for a typed Rust reference kernel;
  \item a structural formalization target for Lean;
  \item a composition boundary at which external primitive security assumptions can be
        attached without being hidden inside protocol lemmas.
\end{enumerate}

The protocol is intentionally narrower than a complete production messaging system. It
does not specify federation, servers, DHT placement, relay topology, account management,
GUI behavior, or global traffic-analysis resistance. These omissions are not accidental:
they keep the application state authority and the cryptographic research question
auditable.

\subsection{Contributions}

This paper makes the following contributions.
\begin{enumerate}[leftmargin=*]
  \item It defines a reproducible protocol core with typed state ownership, strict turn
        semantics, authenticated public receiver packages, and atomic package rotation.
  \item It specifies a canonical wire format with fixed field order, byte lengths, object
        tags, associated-data construction, and a 4096-byte envelope.
  \item It gives a complete adversary model in which one endpoint may be continuously
        controlled while the other endpoint's private state and fresh randomness remain
        outside the adversary's direct input interface.
  \item It separates structural theorems from computational assumptions and defines the
        exact correspondence required between a generic PPT adversary and an adaptive
        protocol-game adversary.
  \item It provides a conditional reduction shape that makes every primitive certificate
        identify the same security parameter, concrete wrapper game, probability model,
        and PPT witness.
  \item It defines reproducibility artifacts: a deterministic test randomness schedule,
        machine-readable vectors, a reference implementation, and an end-to-end fixture.
\end{enumerate}

\subsection{Claim boundary}

The principal claim of this paper is not that the protocol is unconditionally secure in
all implementations. The claim is instead layered:

\begin{enumerate}[leftmargin=*]
  \item the protocol state machine and wire contract are explicit;
  \item the stated structural properties are checked under the premises of the Lean
        model;
  \item a supplied concrete game instance can be related to named primitive wrapper games
        by an explicit hybrid bound;
  \item computational security follows only when the external primitive certificates and
        implementation assumptions required by that bound are supplied.
\end{enumerate}

In particular, the formalization does not internally prove the computational security of
the external X25519, ML-KEM-768, HKDF, ChaCha20-Poly1305, Ed25519, or OS CSPRNG
implementations. This distinction is part of the protocol's trust model.

\section{Background and Position}
\label{sec:background}

Link Chat uses standard cryptographic interfaces rather than proposing a new primitive.
The classical component is an X25519 wrapper with an ephemeral public value as its wire
ciphertext \cite{rfc7748}. The post-quantum component is ML-KEM-768, used through the
standard KEM interface defined by FIPS 203 \cite{fips203}. Key derivation uses HKDF-SHA-256
\cite{rfc5869}; authenticated encryption uses ChaCha20-Poly1305 \cite{rfc8439}; long-term
identity binding uses Ed25519 \cite{rfc8032}.

The protocol should therefore be read as a state-management and composition study. Its
novelty claim is deliberately limited to the arrangement of state authority, turn
ownership, package publication, and proof boundaries. It does not claim a new hardness
assumption, a new KEM, a new AEAD, or a general replacement for mature messaging
protocols.

\section{System Model and Notation}
\label{sec:model}

\subsection{Participants and sessions}

There are two endpoints, denoted Alice and Bob. A session has a fresh identifier
$\mathsf{sid}$, a protocol version $v$, and a cipher-suite identifier $\mathsf{suite}$.
The v1.0 suite is:

\begin{lstlisting}[caption={The frozen v1.0 cipher suite}]
V1HybridMlKem768 = Ed25519 + X25519 wrapper + ML-KEM-768
                   + HKDF-SHA-256 + SHA-256 + ChaCha20-Poly1305
\end{lstlisting}

Each endpoint has a long-term Ed25519 identity key pair and an independent local receiver
secret package. The endpoints never exchange their receiver private keys.

\subsection{Security parameter and probability model}

Let $\lambda$ be the security parameter. A concrete game instance must provide a sample
space $\Omega_\lambda$, all random variables used by the experiment, a probability
measure over $\Omega_\lambda$, an adversary strategy, resource bounds, a success predicate,
and an advantage definition. In the finite Lean model, a typical sample space is a finite
uniform product:
\[
  \Omega_\lambda = \prod_{i=1}^{m(\lambda)} \{0,1\}^{\ell_i(\lambda)},
\]
where each coordinate is consumed by a declared random-source event. This finite model is
an explicit semantics for the checked instance; it is not automatically an OS-randomness
model.

For a bit-valued experiment $G^b_\lambda(A)$, define
\[
  \Pr[G^b_\lambda(A)=1]
\]
as the probability over the specified sample space and experiment coins. The distinguishing
advantage is
\[
  \Adv_G(A,\lambda) = \left|
    \Pr[G^0_\lambda(A)=1] - \Pr[G^1_\lambda(A)=1]
  \right|.
\]

\subsection{PPT adversaries and witnesses}

The formal endpoint used for the audited theorem is a supplied adversary witness
$A_\lambda$ together with a certificate that it is probabilistic polynomial time. The
audited theorem is intentionally an instance theorem:
\[
  \forall \lambda \in \mathbb{N},\quad
  \text{given } A_\lambda \text{ and its PPT witness, the stated bound holds.}
\]
It is not silently upgraded to a universal theorem over every possible external PPT
implementation. Such an upgrade would require a separate semantic universe of adversary
programs and a proof that every program is represented by the protocol-game interface.

\section{Protocol Architecture}
\label{sec:architecture}

\subsection{Layering}

The protocol is divided into four layers:
\[
\text{Identity/Session}
\rightarrow
\text{Receiver State}
\rightarrow
\text{KEM/KDF/Message}
\rightarrow
\text{Transport/Mailbox/Path}.
\]
Each layer exposes an interface to the next layer. A lower-layer event cannot silently
authorize an unrelated application state transition.

\subsection{Endpoint state}

For endpoint $X \in \{A,B\}$, the production state is:
\begin{align*}
  \State_X = \{&\mathsf{SessionContext},
  \mathsf{IK}_{X,\mathrm{sk}},
  \mathsf{IK}_{\mathrm{peer},\mathrm{pk}},\\
  &\mathsf{Package}^{\mathrm{priv}}_X,
  \mathsf{Package}^{\mathrm{pub}}_{\mathrm{peer}},
  \mathsf{expectedTurn},\\
  &\mathsf{consumedMessageIDs},
  \mathsf{packageChainCheckpoint}\}.
\end{align*}

The private package contains the local X25519 and ML-KEM private material. Its public
projection contains no private key, message key, path secret, or CSPRNG state. The public
projection is the only package representation that can be serialized, logged, exported,
or passed to mailbox code.

\subsection{State authority}

The application message can carry authenticated encrypted data and a signed public return
package. It can request one valid state-machine transition. It cannot be interpreted as:

\begin{lstlisting}[caption={Forbidden interpretations of an application message}]
SET_RECEIVER_PRIVATE_KEY
SET_RECEIVER_TOKEN
SET_RECEIVER_GENERATION
ROLLBACK_RECEIVER_PACKAGE
GRANT_STATE_OWNERSHIP
\end{lstlisting}

Only the protocol transition and validated storage recovery may create a new private
receiver package. ACK, SACK, retry, mailbox replication, route failure, and ordinary
filesystem events have no application-state authority.

\section{Identity and Session Initialization}
\label{sec:initialization}

\subsection{Out-of-band binding}

The initial procedure must authenticate at least:
\[
  (v,\mathsf{suite},\mathsf{IK}_{A,\mathrm{pk}},\mathsf{IK}_{B,\mathrm{pk}},
    \mathsf{sid},S_0,\mathsf{Package}^{\mathrm{pub}}_A,
    \mathsf{Package}^{\mathrm{pub}}_B,
    \mathsf{Root}_A,\mathsf{Root}_B).
\]
Identity fingerprints must be confirmed through a physical or equivalent authenticated
channel. A network-delivered first identity claim is not sufficient.

The canonical identity-binding input is:
\[
  \mathsf{bind} = \mathsf{Canon}(
  \texttt{LinkChat/identity-binding/v1},v,\mathsf{suite},\mathsf{sid},
  \mathsf{IK}_{A,\mathrm{pk}},\mathsf{IK}_{B,\mathrm{pk}}).
\]
Both endpoints sign the same binding under their own identity keys. An identity change
invalidates the session and requires a new session identifier, bootstrap secret, receiver
state, package chain, and mailbox-token set.

\subsection{Bootstrap}

Let $S_0$ be the out-of-band bootstrap secret. Define:
\begin{align*}
  \mathsf{PRK}_{\mathrm{boot}} &= \mathsf{HKDF\mbox{-}Extract}(0^{256},S_0),\\
  K_{\mathrm{boot}} &= \mathsf{HKDF\mbox{-}Expand}(\mathsf{PRK}_{\mathrm{boot}},
       \mathsf{Canon}(\texttt{LinkChat/bootstrap/v1},v,\mathsf{suite},\mathsf{sid}),32).
\end{align*}
The bootstrap values are not the direct root of later receiver private states. After
bootstrap confirmation, $S_0$, $\mathsf{PRK}_{\mathrm{boot}}$, and $K_{\mathrm{boot}}$
enter the implementation's erase procedure.

\subsection{Initial receiver packages}

Each endpoint independently samples:
\[
\begin{aligned}
  (xsk_X,xpk_X) &\sample \mathsf{X25519.KeyGen}(\mathsf{CSPRNG}),\\
  (psk_X,ppk_X) &\sample \mathsf{MLKEM768.KeyGen}(\mathsf{CSPRNG}),\\
  \mathsf{Token}_X &\sample \{0,1\}^{256}.
\end{aligned}
\]
The initial public package also contains a locally unique KeyID, expiration information,
the package-chain root predecessor, and an Ed25519 authentication value.

For owner $X$, the chain root is:
\[
  \mathsf{Root}_X = \mathsf{SHA256}(\mathsf{Canon}(
  \texttt{LinkChat/package-chain-root/v1},v,\mathsf{suite},\mathsf{sid},
  \mathsf{IK}_{A,\mathrm{pk}},\mathsf{IK}_{B,\mathrm{pk}},X)).
\]

\section{Cryptographic Construction}
\label{sec:crypto}

\subsection{Primitive profile}

\begin{table}[ht]
\centering
\caption{Frozen v1.0 primitive profile.}
\label{tab:primitives}
\begin{tabularx}{\textwidth}{P{0.28\textwidth}Y}
\toprule
Function & Primitive and parameters \\
\midrule
Identity signature & Ed25519, 32-byte public key, 64-byte signature \\
Classical wrapper & X25519 ephemeral DH, 32-byte wire ciphertext \\
Post-quantum KEM & ML-KEM-768, 1184-byte public key, 1088-byte ciphertext \\
KDF & HKDF-SHA-256, 32-byte derived keys \\
Hash & SHA-256 \\
AEAD & ChaCha20-Poly1305, 12-byte nonce, 16-byte tag \\
Freshness & OS CSPRNG or an equivalent vetted source \\
\bottomrule
\end{tabularx}
\end{table}

\subsection{X25519 wrapper}

The X25519 wrapper represents an ephemeral public value as its wire ciphertext:
\[
  (esk_X,epk_X) \sample \mathsf{X25519.KeyGen}(\mathsf{CSPRNG}),\qquad
  ct_X=epk_X,\qquad
  ss_X=\mathsf{X25519}(esk_X,pk_{X,\mathrm{recv}}).
\]
The receiver computes the same shared value using its private scalar and $ct_X$. Low-order
inputs and an all-zero shared result are rejected. The wrapper's exact wire ciphertext
length is 32 bytes.

\subsection{ML-KEM-768}

The post-quantum component uses:
\[
  \begin{aligned}
    (pk_{PQ},sk_{PQ}) &\sample \mathsf{MLKEM768.KeyGen}(\mathsf{CSPRNG}),\\
    (ct_{PQ},ss_{PQ}) &\sample \mathsf{MLKEM768.Encapsulate}(pk_{PQ,\mathrm{recv}}).
  \end{aligned}
\]
The receiver decapsulates $ct_{PQ}$ with its local private key. Invalid ciphertexts are
mapped to a stable cryptographic failure category and must not expose backend-specific
failure details through the remote interface.

\subsection{Hybrid derivation}

The two shared secrets are combined only after domain-separated framing:
\[
  \mathsf{HybridInput}=\mathsf{Canon}(\texttt{LinkChat/hybrid/v1},v,\mathsf{suite},
      \mathsf{sid},\mathsf{turn},\mathsf{receiverKeyID},ct_X,ct_{PQ},ss_X,ss_{PQ}).
\]
Then:
\[
  \mathsf{PRK}_{\mathrm{hybrid}}=
  \mathsf{HKDF\mbox{-}Extract}(0^{256},\mathsf{HybridInput}).
\]
Neither KEM shared secret is used directly as an AEAD key.

\subsection{Message and header keys}

Define:
\begin{align*}
  \mathsf{MessageInfo} &= \mathsf{Canon}(\texttt{LinkChat/message/v1},v,\mathsf{suite},
    \mathsf{sid},\mathsf{turn},\mathsf{direction},\mathsf{messageID},\mathsf{receiverKeyID}),\\
  \mathsf{HeaderInfo} &= \mathsf{Canon}(\texttt{LinkChat/header/v1},v,\mathsf{suite},
    \mathsf{sid},\mathsf{turn},\mathsf{direction},\mathsf{receiverKeyID}),\\
  MK &= \mathsf{HKDF\mbox{-}Expand}(\mathsf{PRK}_{\mathrm{hybrid}},\mathsf{MessageInfo},32),\\
  HK &= \mathsf{HKDF\mbox{-}Expand}(\mathsf{PRK}_{\mathrm{hybrid}},\mathsf{HeaderInfo},32).
\end{align*}
The two keys are distinct by domain separation. A Message Key is used for one application
message only.

\subsection{Nonces}

The message nonce is derived from a domain-separated nonce input:
\[
  \mathsf{NonceInfo}=\mathsf{Canon}(\texttt{LinkChat/nonce/v1},\mathsf{sid},
     \mathsf{turn},\mathsf{direction},\mathsf{messageID},\mathsf{receiverKeyID}),
\]
\[
  n_M=\mathsf{HKDF\mbox{-}Expand}(\mathsf{HKDF\mbox{-}Extract}(0^{256},MK),
       \mathsf{NonceInfo},12).
\]
Header nonce derivation uses the same context shape under an independent header-nonce
domain. Nonce uniqueness depends on correct Turn,
Direction, MessageID, KeyID, and Message Key binding.

\subsection{Package hash and authentication}

Let $\mathsf{PackageBody}$ be the canonical package without $\mathsf{packageAuth}$. Then:
\[
  \mathsf{packageHash}=\mathsf{SHA256}(\mathsf{Canon}(
    \texttt{LinkChat/package-hash/v1},\mathsf{PackageBody})),
\]
\[
  \begin{aligned}
    \mathsf{packageAuth}={}&\mathsf{Ed25519.Sign}(IK_{X,\mathrm{sk}},\\
      &\mathsf{Canon}(\texttt{LinkChat/receiver-package-auth/v1},
        \mathsf{PackageBody})).
  \end{aligned}
\]
The verifier recomputes the hash, verifies the signature under the pinned identity key,
and checks the chain predecessor.

\subsection{AEAD and associated data}

The logical Header is sealed as:
\[
  C_H=\mathsf{Seal}(HK,n_H,\mathsf{HeaderAD},\mathsf{Canon}(\mathsf{Header})),
\]
and the fixed PayloadFrame is sealed as:
\[
  C_M=\mathsf{Seal}(MK,n_M,\mathsf{MessageAD},\mathsf{PayloadFrame}).
\]

The HeaderAD is:
\[
  \begin{aligned}
    \mathsf{HeaderAD}=\mathsf{Canon}(&\texttt{LinkChat/header-ad/v1},v,\mathsf{suite},\\
      &\mathsf{sid},\mathsf{turn},\mathsf{direction},\mathsf{receiverKeyID}).
  \end{aligned}
\]
The encrypted Header contains MessageID, so MessageID cannot appear in HeaderAD. After
successful Header open and canonical decoding, the receiver constructs:
\[
  \begin{aligned}
    \mathsf{MessageAD}=\mathsf{Canon}(&\texttt{LinkChat/message-ad/v1},v,\mathsf{suite},\\
      &\mathsf{sid},\mathsf{turn},\mathsf{direction},\mathsf{messageID},
        \mathsf{receiverKeyID}).
  \end{aligned}
\]
AEAD open is a single decrypt-and-authenticate operation. Failure stops processing and
cannot update consumed IDs or generate a new package.

\section{Canonical Wire Format}
\label{sec:wire}

\subsection{Object encoding}

Every canonical object begins with a nine-byte header:
\begin{center}
\begin{tabular}{ccl}
\toprule
Offset & Size & Meaning \\
\midrule
0 & 1 & Object tag \\
1 & 2 & Canonical format version, unsigned big-endian \\
3 & 2 & Protocol major version, unsigned big-endian \\
5 & 2 & Protocol minor version, unsigned big-endian \\
7 & 2 & Field count, unsigned big-endian \\
\bottomrule
\end{tabular}
\end{center}

Every field is encoded as $\mathsf{u32}_{\mathrm{be}}(\mathsf{length})$ followed by the
field bytes. Fields are neither omitted nor reordered. Decoders reject truncation,
overflow, unsupported tags or versions, invalid field counts, semantic length errors,
trailing bytes, and non-canonical re-encoding.

\subsection{Object registry}

\begin{table}[ht]
\centering
\caption{v1.0 object tags and field counts.}
\label{tab:objects}
\begin{tabular}{clc}
\toprule
Tag & Object & Fields after object header \\
\midrule
0x01 & ApplicationMessage & 8 \\
0x02 & ReceiverPackage & 11 \\
0x03 & Envelope & 6 \\
0x04 & ACK & 6 \\
0x05 & SACK & 6 \\
0x06 & MailboxRecord & 2 \\
0x07 & Header & 10 \\
0x7f & Context & variable \\
\bottomrule
\end{tabular}
\end{table}

\subsection{ReceiverPackage}

The public package has eleven fields:
\begin{center}
\begin{tabular}{cll}
\toprule
Index & Field & Encoding \\
\midrule
0 & cipher suite & u16 \\
1 & session ID & u64 \\
2 & turn & u64 \\
3 & generation & u64 \\
4 & key ID & u64 \\
5 & X25519 public key & 32 bytes \\
6 & ML-KEM-768 public key & 1184 bytes \\
7 & mailbox token & 32 bytes \\
8 & expiration & u64 \\
9 & previous package hash & 32 bytes \\
10 & package authentication & 64 bytes \\
\bottomrule
\end{tabular}
\end{center}
Its canonical encoded length is exactly 1439 bytes. Private X25519 and ML-KEM values are
not fields of this object.

\subsection{Header and Envelope}

The ten Header fields are:
\[
  \begin{aligned}
    (&\mathsf{suite},\mathsf{sid},\mathsf{turn},\mathsf{direction},\mathsf{messageID},\\
     &\mathsf{receiverKeyID},\mathsf{messageType},\mathsf{senderPackageHash},\\
     &\mathsf{receiverPackageHash},\mathsf{returnReceiverPackage}).
  \end{aligned}
\]
The canonical logical Header is 1588 bytes. Header AEAD adds a 16-byte tag, giving an
encrypted Header field of 1604 bytes.

The six Envelope fields are:
\[
  (\mathsf{mailboxToken},ct_X,ct_{PQ},C_H,C_M,\mathsf{padding}).
\]
Their fixed offsets are:
\begin{center}
\begin{tabular}{lrr}
\toprule
Field & Offset range & Value length \\
\midrule
Object header & 0--8 & 9 \\
Mailbox token & 9--44 & 32 \\
X25519 ciphertext & 45--80 & 32 \\
ML-KEM ciphertext & 81--1172 & 1088 \\
Encrypted Header & 1173--2780 & 1604 \\
Message ciphertext & 2781--4091 & 1307 \\
Padding & 4092--4095 & 0 \\
\bottomrule
\end{tabular}
\end{center}
The complete canonical Envelope is exactly 4096 bytes.

\subsection{PayloadFrame}

For application bytes $m$ of length $\ell \leq 1289$:
\[
  \mathsf{PayloadFrame}(m)=\mathsf{u16}_{\mathrm{be}}(\ell)\mathbin{||}m\mathbin{||}0^{1289-\ell}.
\]
The frame is exactly 1291 bytes. ChaCha20-Poly1305 adds a 16-byte tag, producing a
1307-byte message ciphertext.

\subsection{ACK, SACK, and MailboxRecord}

ACK and SACK are transport/control objects. Their protocol version is carried by the
common object header. ACK has payload fields $(\mathsf{suite},\mathsf{sid},\mathsf{turn},
\mathsf{direction},\mathsf{messageID},\mathsf{deliveryStatus})$. SACK has payload fields
$(\mathsf{suite},\mathsf{sid},\mathsf{turn},\mathsf{direction},\mathsf{baseMessageID},
\mathsf{bitmap})$. A MailboxRecord binds one opaque token to one canonical Envelope.

None of these control objects can advance a turn, rotate a receiver package, mark an
application message consumed, or restore an erased message key.

\section{Endpoint Algorithms}
\label{sec:algorithms}

\subsection{Send preparation}

On an expected turn $t$, the endpoint first checks that it is the leader:
\[
  \mathsf{Leader}(t)=
  \begin{cases}
    \mathsf{Alice}, & t \bmod 2=0,\\
    \mathsf{Bob}, & t \bmod 2=1.
  \end{cases}
\]
The send procedure then selects the current verified peer package, samples ephemeral
X25519 material, encapsulates under both receiver public keys, derives the hybrid keys,
constructs the Header and PayloadFrame, seals both, and returns an Envelope plus a
complete pending state. Preparation alone does not publish state.

\begin{lstlisting}[caption={Abstract send preparation}]
prepare_send(plaintext):
    require local_endpoint == Leader(expected_turn)
    require peer_package is current, verified, unexpired
    sample ephemeral X25519 key
    encapsulate to peer X25519 and ML-KEM public keys
    derive PRK_hybrid, HK, MK, header_nonce, message_nonce
    construct Header with current sender package
    construct fixed PayloadFrame(plaintext)
    seal Header with HeaderAD
    seal PayloadFrame with MessageAD
    construct 4096-byte Envelope
    return PreparedSend(envelope, candidate_state)
\end{lstlisting}

\subsection{Receive processing}

The receive order is normative:
\begin{lstlisting}[caption={Normative receive processing order}]
bounded decode of the 4096-byte Envelope
outer Token and context checks
X25519 and ML-KEM decapsulation
Hybrid KDF
Header AEAD Open with HeaderAD
canonical Header decode
Session, Turn, Direction, MessageID, and KeyID checks
receiver package hash check
return Package signature and chain checks
Message KDF
Message AEAD Open with MessageAD
PayloadFrame length and zero-padding checks
duplicate and consumed-message checks
fresh local Receiver Package generation
complete pending-state construction
durable prepare and atomic commit
application delivery
\end{lstlisting}

The order prevents an unauthenticated field from becoming a state-authority input. In
particular, the receiver does not generate and persist a new package until the entire
message, authenticated return package, and payload have passed validation.

\subsection{Successful receive transition}

After all checks succeed, the receiver samples a new local package:
\[
\begin{aligned}
  (xsk',xpk') &\sample \mathsf{X25519.KeyGen}(\mathsf{CSPRNG}),\\
  (psk',ppk') &\sample \mathsf{MLKEM768.KeyGen}(\mathsf{CSPRNG}),\\
  \mathsf{Token}' &\sample \{0,1\}^{256},\\
  \mathsf{generation}' &= \mathsf{generation}+1,\\
  \mathsf{previousHash}' &= \mathsf{Hash}(\mathsf{currentPublicPackage}).
\end{aligned}
\]
The new private material, public projection, expected turn, peer package, and consumed
MessageID set are committed as one logical state.

\subsection{Failure behavior}

Every following condition is a no-state-change result:
\begin{itemize}[leftmargin=*]
  \item malformed, non-canonical, truncated, oversized, or unsupported input;
  \item session, token, turn, direction, key-ID, expiration, or context mismatch;
  \item replay, future turn, duplicate message, package-hash, signature, or chain failure;
  \item X25519, ML-KEM, Header AEAD, Message AEAD, or PayloadFrame failure;
  \item storage preparation, durability, or commit failure.
\end{itemize}
No failure changes the current package, peer package, expected turn, consumed IDs,
generation, or retained message-key state. A rejected plaintext is not delivered.

\section{Storage and Recovery}
\label{sec:storage}

\subsection{Prepare/commit contract}

The abstract storage interface is:
\[
  \begin{aligned}
    \mathsf{load}() &\rightarrow \mathsf{CurrentState}, &
    \mathsf{prepare\_commit}(p) &\rightarrow \mathsf{PreparedCommit},\\
    \mathsf{commit}(c) &\rightarrow \mathsf{Result}, &
    \mathsf{recover}() &\rightarrow \mathsf{CurrentState}.
  \end{aligned}
\]

A recommended durable order is:
\begin{enumerate}[leftmargin=*]
  \item write and validate Pending;
  \item complete a durability barrier;
  \item write the commit marker;
  \item complete a durability barrier;
  \item atomically publish Current;
  \item complete the required directory or medium barrier;
  \item retire the old state;
  \item attempt implementation-defined zeroization.
\end{enumerate}
The old private state must not be erased before the new state is durable.

\subsection{Recovery matrix}

\begin{table}[ht]
\centering
\caption{Required recovery behavior.}
\label{tab:recovery}
\begin{tabularx}{\textwidth}{P{0.42\textwidth}Y}
\toprule
Crash or corruption condition & Required result \\
\midrule
Pending incomplete & Discard Pending; retain Current \\
Pending complete without marker & Discard Pending; retain Current \\
Marker present without published Current & Adopt Pending as Current \\
Current published while old state remains & Use new Current; retire old later \\
Record hash mismatch & Refuse normal recovery and report storage corruption \\
Package-chain mismatch & Refuse normal recovery and report rollback or fork \\
Private/public projection mismatch & Refuse recovery \\
\bottomrule
\end{tabularx}
\end{table}

Atomic rename, filesystem synchronization, or ordinary deletion does not by itself prove
rollback resistance or physical erasure. A true rollback-resistance claim requires an
external monotonic property, such as an appropriate hardware-backed or service-backed
counter.

\section{Security Games}
\label{sec:games}

\subsection{Complete adversary view}

In the Alice-compromised experiment, the adversary may choose or observe Alice's local
state, plaintexts, randomness, ciphertexts, tokens, application actions, network actions,
transport actions, mailbox and DHT actions, and replay, delay, drop, modification, and
injection behavior. It also receives Bob's declared public outputs and protocol-level
observations.

The complete view is modeled as:
\begin{align*}
  \View_A = (&\mathsf{controlledState},\mathsf{plaintexts},\mathsf{randomness},
  \mathsf{ciphertexts},\mathsf{tokens},\\
  &\mathsf{networkTrace},\mathsf{transportTrace},\mathsf{mailboxTrace},
  \mathsf{DHTTrace},\\
  &\mathsf{peerPublicOutputs},\mathsf{ACKError},\mathsf{acceptReject},
  \mathsf{eventTime},\\
  &\mathsf{logicalLength},\mathsf{envelopeLength},
  \mathsf{stateUpdateCount},\mathsf{packageUpdateCount}).
\end{align*}
If an experiment exposes fewer components, it is a projection rather than the complete
view.

\subsection{Structural state non-interference}

Let $s$ be an honest state, $\rho$ a fixed local fresh schedule, and $x_1,x_2$ peer inputs.
The strong structural property is:
\begin{definition}[State non-interference]
If both inputs satisfy the same acceptance branch, then
\[
  \mathsf{nextSecret}(s,x_1,\rho)=\mathsf{nextSecret}(s,x_2,\rho).
\]
\end{definition}
The property is stronger than merely stating that peer bytes do not appear syntactically
in a state field. It says that, for a fixed local fresh schedule, peer choice cannot choose
the honest endpoint's next private secret.

\subsection{Future-state challenge}

The Real world uses the honest endpoint's generated future receiver state. The Random world
uses an independent random future-state value. Both worlds use the same initial state,
adaptive strategy, trace schedule, timing, length metadata, ACK/error interface,
accept/reject interface, and challenge message inputs. Only the challenge material derived
from the future state may differ.

\subsection{Future-message challenge}

For two equal-length messages $M_0,M_1$, sample $b\sample\{0,1\}$ and return:
\[
  C^*=\mathsf{Seal}(K_b,n,AD,\mathsf{PayloadFrame}(M_b)).
\]
The challenge is a real-format ciphertext inside a fixed-size envelope, not an exposed
internal state value. The game therefore tests message indistinguishability through the
same observable interface used by the protocol.

\subsection{FS, PCS, authentication, and KCI}

Forward secrecy gives the adversary a current state and challenges it on a previously
erased message key or message. Post-compromise security gives the adversary state at time
$t_0$, restores honest behavior at $t_1$, requires a fresh local package, and challenges a
message after recovery. Authentication challenges package forgery across sessions, turns,
KeyIDs, or package chains. KCI is split into the two directions of identity-key
compromise; an affected session must be invalidated and re-established through identity
binding.

\section{Structural Properties and Formalization}
\label{sec:formalization}

The associated Lean project is organized into protocol definitions, structural games,
standard-crypto interfaces, parameterized probability semantics, concrete reduction
bookkeeping, implementation semantics, and trust-boundary audits. The relevant conceptual
modules include:
\begin{lstlisting}[caption={Conceptual organization of the Lean development}]
LinkChat.lean
SecurityGames.lean
RefinedProtocol.lean
StandardCrypto.lean
ComputationalGames.lean
ParameterizedConcreteSecurity.lean
ConcreteSECCReduction.lean
ConcreteSecurityReductions.lean
ImplementationSemantics.lean
TrustBoundaryAudit.lean
AuditHardening.lean
\end{lstlisting}

Under their explicit premises, the formalized structural layer covers:
\begin{itemize}[leftmargin=*]
  \item strict alternation and unique leader selection;
  \item invalid-input no-state-change behavior;
  \item replay, future-turn, duplicate, and rollback handling;
  \item atomic receiver-package rotation;
  \item package binding and state-injection resistance;
  \item bidirectional state non-interference;
  \item continuous adaptive attack traces;
  \item abstract KEM and key-derivation correctness.
\end{itemize}

The formalization also provides interfaces for finite and parameterized probability
semantics, PPT resource certificates, the complete adversary view, future-state and
future-message challenges, FS/PCS/authentication/KCI games, named primitive wrapper games,
hybrid advantage addition, and finite-term negligible closure.

\subsection{Adversary refinement}

The generic PPT adversary must be related to the protocol-game adversary rather than merely
given the same name. The required refinement theorem has the form:
\[
  \forall h,\quad
  \mathsf{protocolAdversary}(A,h)
  =\mathsf{refine}(A,\mathsf{CompleteAdversaryView}(h)),
\]
where $h$ is the complete observable history and the refinement preserves the action,
resource, and observation interfaces. This establishes that the protocol game can make
adaptive decisions from the same history that the abstract PPT witness is allowed to see.

The audited theorem additionally requires that the challenge, event time, logical length,
fixed envelope length, ACK/error values, accept/reject observations, state-update count,
and package-update count are the same components on both sides of the binding.

\section{Conditional Computational Theorem}
\label{sec:computational}

\subsection{Primitive certificates}

Each certificate must identify all of the following:
\begin{enumerate}[leftmargin=*]
  \item the designated wrapper game, not merely a primitive name;
  \item the same probability model used by the protocol game;
  \item the same security parameter $\lambda$ or index $n$;
  \item a declared PPT adversary witness and its resource bound;
  \item the simulator and interface used by the corresponding hybrid hop.
\end{enumerate}

Thus, a field labelled ``x25519'' is not enough. A valid certificate is a statement about
the exact X25519 wrapper game invoked by the concrete reduction.

\subsection{Hybrid bound}

For a supplied concrete protocol game instance and witnesses, the intended reduction is:
\begin{align}
  \Adv_{\mathrm{SECC}}(A,\lambda)
  \leq{}&\Adv_{\mathrm{State}}(A,\lambda)
  +\Adv_{\mathrm{X25519}}(B_X,\lambda)
  +\Adv_{\mathrm{MLKEM}}(B_{PQ},\lambda) \notag\\
  &+\Adv_{\mathrm{HKDF}}(B_{KDF},\lambda)
  +\Adv_{\mathrm{AEAD,conf}}(B_C,\lambda) \notag\\
  &+\Adv_{\mathrm{AEAD,int}}(B_I,\lambda)
  +\Adv_{\mathrm{CSPRNG}}(B_R,\lambda).
  \label{eq:mainbound}
\end{align}
The state term is zero when Real and State-Isolated are extensionally equal under the
same observable interface. Otherwise it remains an explicit term.

\begin{theorem}[Audited concrete-instance SECC theorem]
\label{thm:secc}
Let $G^0_\lambda$ and $G^1_\lambda$ be one concrete Real/Random protocol-game pair with
the same complete observable interface. Let $A_\lambda$ be a supplied PPT witness with an
adversary-refinement theorem. Suppose the concrete reduction supplies wrapper-game
certificates for X25519, ML-KEM-768, HKDF-SHA-256, ChaCha20-Poly1305 confidentiality and
integrity, and the honest fresh-randomness source, all at the same parameter $\lambda$.
If the hybrid hops satisfy \cref{eq:mainbound}, then:
\[
  \Adv_{\mathrm{SECC}}(A_\lambda,\lambda)
  \leq \text{the right-hand side of Equation~\eqref{eq:mainbound}}.
\]
If, in addition, the state term is zero and every primitive term is negligible in
$\lambda$, then:
\[
  \Adv_{\mathrm{SECC}}(A_\lambda,\lambda)\leq \negl(\lambda).
\]
\end{theorem}

\begin{proof}
The first inequality is the triangle inequality applied to the explicitly declared game
sequence. Each adjacent hybrid changes one named component and is bounded by its wrapper
certificate. Adversary refinement allows the same adaptive strategy to be replayed in the
simulator because the simulator receives the same complete history and exposes the same
challenge and metadata. If Real and State-Isolated are extensionally equal, their event
probabilities are equal and the state term is zero. Finally, a finite sum of negligible
functions is negligible, so the additional premises imply the final conclusion.
\end{proof}

\begin{remark}
The theorem is deliberately not a theorem that every PPT adversary, every implementation,
or every external cryptographic library is secure. It is a theorem for the supplied,
audited, parameterized game instance and its declared witnesses.
\end{remark}

\subsection{What the theorem does not establish}

The theorem does not by itself establish physical erasure, rollback resistance on an
ordinary filesystem, resistance to memory disclosure, resistance to side channels, secure
cross-process coordination, availability against a malicious mailbox, anonymity, or
traffic-analysis protection. Those are separate implementation or deployment properties.

\section{Reference Implementation and Reproducibility}
\label{sec:reproducibility}

The Rust reference kernel is divided into typed, maintainable crates:
\begin{center}
\begin{tabular}{ll}
\toprule
Crate & Role \\
\midrule
\code{linkchat-types} & fixed-size semantic identifiers and public types \\
\code{linkchat-protocol} & canonical objects, codecs, contexts, and wire errors \\
\code{linkchat-crypto} & typed primitive-provider boundary \\
\code{linkchat-core} & pure protocol transitions and package logic \\
\code{linkchat-engine} & prepare/commit endpoint orchestration \\
\code{linkchat-storage} & durable transaction and recovery adapters \\
\code{linkchat-transport} & packet, retry, ACK, SACK, and metadata boundary \\
\code{linkchat-mailbox} & token-bound mailbox interface and in-memory adapter \\
\code{linkchat-verification} & properties, vectors, fuzz smoke, and differential checks \\
\code{linkchat-ffi} & opaque-handle C ABI boundary \\
\code{linkchat-cli} & inspection and vector commands \\
\bottomrule
\end{tabular}
\end{center}

For deterministic test artifacts, define:
\[
  R(\mathsf{seed},i)=\mathsf{SHA256}(\mathsf{seed}\mathbin{||}\mathsf{u64}_{\mathrm{be}}(i)).
\]
This schedule is test-only. Production implementations must use an operating-system CSPRNG
or an equivalent vetted source.

The published vector set includes identity binding, bootstrap, package structure,
application-message encoding, context AAD, payload frame, envelope layout, HKDF framing,
AEAD round trips, KEM boundaries, and replay/recovery behavior. Structural vectors use
patterned bytes and exact lengths; cryptographic test material is explicitly marked
test-only.

The reference kernel was checked with:
\begin{lstlisting}[caption={Reference verification commands}]
cargo fmt --all -- --check
cargo check --workspace
cargo test --workspace
cargo clippy --workspace --all-targets --all-features -- -D warnings
cargo run -p linkchat-cli -- test-vector
\end{lstlisting}
These commands provide implementation and vector evidence. They do not replace the Lean
theorems or external primitive-security certificates.

\subsection{Turn-zero fixture}

A complete reproducible fixture uses $\mathsf{sid}=42$, Turn $0$, Alice as leader,
$\mathsf{messageID}=7$, receiver KeyID $99$, expiration $1700000000$, and UTF-8 payload
\texttt{hello}. The PayloadFrame is:
\[
  \texttt{0005}\mathbin{||}\texttt{68656c6c6f}\mathbin{||}0^{1284},
\]
with length 1291 bytes. The resulting Envelope has length 4096 bytes and contains no
plaintext message ID, plaintext payload, private key, or secret state.

\section{Discussion}
\label{sec:discussion}

The protocol's main design choice is to separate input authentication from state authority.
An application message can be accepted, rejected, replayed, or delivered, but it cannot
directly select the next private receiver state. This produces a clean structural property
that is useful before any computational assumption is applied.

The strict alternating turn rule is a deliberate restriction. It avoids ambiguity from
concurrent updates, generic receive windows, and multiple valid commits for one endpoint.
Transport may perform retries or parallel mailbox retrieval, but application commit is
serialized. A future extension that introduces windows, arbitrary reordering, or concurrent
commits must revise the state model, wire semantics, security games, and proof obligations.

The fixed envelope and payload frame make the protocol reproducible and simplify comparison
of observable lengths. They do not, however, constitute a global traffic-analysis defense.
The adversary model may expose logical length, event timing, ACK/error results, and state
update counts precisely because the goal is to avoid hiding these channels in an informal
security claim.

\section{Limitations and Research Boundary}
\label{sec:limitations}

The following boundaries are explicit.
\begin{enumerate}[leftmargin=*]
  \item \textbf{Primitive security:} The paper assumes named computational security
        certificates for the concrete wrappers. It does not prove the low-level security
        of X25519, ML-KEM-768, HKDF, AEAD, Ed25519, or the OS CSPRNG inside Lean.
  \item \textbf{Storage:} The protocol specifies a crash-safe transaction contract, but
        rollback resistance requires an external monotonic property and physical erasure
        remains platform-dependent.
  \item \textbf{Deployment:} Servers, DHTs, relays, production mailboxes, account
        systems, and network availability are outside the core.
  \item \textbf{Side channels:} Memory disclosure, timing leakage, cache behavior,
        swap, crash dumps, backups, and allocator copies are not modeled as eliminated.
  \item \textbf{Adversary theorem scope:} The audited final theorem is for a supplied
        parameterized game instance and PPT witness with an explicit semantic refinement.
        It is not an automatic universal quantification over arbitrary programs.
  \item \textbf{Privacy scope:} The core does not establish anonymity, token unlinkability,
        global traffic-analysis resistance, or mailbox availability.
\end{enumerate}

These limitations are part of the research result. They prevent a structural protocol proof
from being mistaken for a complete deployment security certification.

\section{Conclusion}
\label{sec:conclusion}

Link Chat is a reproducible protocol core, based on the Link Chat v1.0 research
specification \cite{linkchat-spec}, for studying how authenticated application
inputs interact with endpoint-local secret-state evolution. Its design uses standard
cryptographic primitives and concentrates the protocol contribution in four mechanisms:
independent receiver state, strict turns, authenticated public packages, and atomic local
rotation after complete validation.

The resulting specification is sufficiently concrete for independent implementation:
object tags, field order, offsets, lengths, cryptographic domains, failure behavior,
storage boundaries, and test vectors are fixed. The associated formalization provides a
structural verification layer and a conditional computational composition layer. The final
computational statement is intentionally honest about its trust boundary: it applies to a
declared concrete game instance with PPT witnesses and named primitive certificates, and it
does not claim that Lean has proved the external libraries or a full messaging deployment.

\appendix

\section{Wire-Length Derivations}
\label{app:lengths}

Each field contributes four bytes of length prefix plus its value length. Thus the public
ReceiverPackage length is:
\[
  9 + 4(11) + 2+8+8+8+8+32+1184+32+8+32+64 = 1439.
\]
The logical Header length is:
\[
  9+4(10)+2+8+8+1+8+8+1+32+32+1439=1588.
\]
The Envelope length is:
\[
  9+4(6)+32+32+1088+1604+1307+0=4096.
\]
The PayloadFrame and ciphertext lengths are:
\[
  2+1289=1291,\qquad 1291+16=1307.
\]

\section{Failure Categories}
\label{app:errors}

The stable remote-visible categories are:
\begin{center}
\begin{tabular}{llll}
\toprule
Malformed & UnsupportedVersion & UnsupportedSuite & ContextMismatch \\
Replay & FutureTurn & Duplicate & PackageInvalid \\
KemFailure & AuthenticationFailure & PayloadInvalid & StorageFailure \\
Unavailable & & & \\
\bottomrule
\end{tabular}
\end{center}
Internal diagnostics may be more detailed, but must not expose private keys, message keys,
complete plaintexts, or CSPRNG state.

\section{Artifact Locations}
\label{app:artifacts}

The English protocol specification and vectors are maintained in the companion
research-document repository:
\begin{lstlisting}
LinkChatDocuments repository root
\end{lstlisting}
The Lean formalization and Rust reference kernel are maintained together in
the companion implementation repository:
\begin{lstlisting}
LinkChat repository root
\end{lstlisting}

\begin{thebibliography}{9}

\bibitem{linkchat-spec}
Link Chat Research Project.
\newblock Link Chat Protocol Research Specification, v1.0-research.
\newblock English protocol specification and reproducibility artifacts, 2026.

\bibitem{rfc7748}
Adam Langley, Mike Hamburg, and Sean Turner.
\newblock Elliptic Curves for Security.
\newblock RFC 7748, Internet Engineering Task Force, 2016.

\bibitem{fips203}
National Institute of Standards and Technology.
\newblock Module-Lattice-Based Key-Encapsulation Mechanism Standard.
\newblock FIPS PUB 203, 2024.

\bibitem{rfc5869}
Hugo Krawczyk and Pasi Eronen.
\newblock HMAC-based Extract-and-Expand Key Derivation Function (HKDF).
\newblock RFC 5869, Internet Engineering Task Force, 2010.

\bibitem{rfc8439}
Yoav Nir and Adam Langley.
\newblock ChaCha20 and Poly1305 for IETF Protocols.
\newblock RFC 8439, Internet Engineering Task Force, 2018.

\bibitem{rfc8032}
Simon Josefsson and Ilari Liusvaara.
\newblock Edwards-Curve Digital Signature Algorithm (EdDSA).
\newblock RFC 8032, Internet Engineering Task Force, 2017.

\bibitem{gamehopping}
Victor Shoup.
\newblock Sequences of Games: A Tool for Taming Complexity in Security Proofs.
\newblock Cryptology ePrint Archive, Report 2004/332, 2004.

\end{thebibliography}

\end{document}
